\documentclass[12pt,a4paper]{article}

\usepackage{geometry}
\geometry{margin=2.5cm}

\usepackage{xcolor}
\usepackage{graphicx}
\usepackage{booktabs}
\usepackage{longtable}
\usepackage{array}
\usepackage{hyperref}

\usepackage{xepersian}

\settextfont{XB Niloofar}
\setlatintextfont{Times New Roman}

\title{
\vspace{-2cm}
\includegraphics[width=3cm]{example-image}\\[1cm]
گزارش اولیه پروژه\\
سامانه MLOps پردازش و پیش‌بینی داده‌های بورس ایران\\
مبتنی بر Apache Airflow
}

\author{
اعضای گروه:\\
متخصص علوم کامپیوتر\\
متخصص علم داده\\
مهندس نرم‌افزار
}

\date{\today}

\begin{document}

\maketitle

\tableofcontents

\newpage

%===================================================

\section{مقدمه}

با رشد سریع داده‌های مالی و توسعه سیستم‌های تحلیل بازار سرمایه، نیاز به طراحی سامانه‌های هوشمند جهت پردازش، مدیریت و تحلیل داده‌های بورسی بیش از پیش احساس می‌شود. بازار بورس ایران نیز روزانه حجم بسیار زیادی از داده‌های مالی، قیمت سهام، شاخص‌ها و اطلاعات معاملاتی تولید می‌کند.

هدف این پروژه طراحی و پیاده‌سازی یک سامانه مبتنی بر \lr{MLOps} برای جمع‌آوری، پردازش، آموزش مدل‌های یادگیری ماشین و ارائه سرویس پیش‌بینی بر پایه داده‌های بازار بورس ایران است.

این پروژه در قالب درس \lr{ML Ops and Data Processing with Apache Airflow} تعریف شده و به‌صورت مرحله‌ای در طول ترم توسعه خواهد یافت.

%===================================================

\section{هدف پروژه}

اهداف اصلی پروژه عبارت‌اند از:

\begin{itemize}
\item جمع‌آوری خودکار داده‌های بازار بورس ایران
\item پردازش و پاکسازی داده‌ها
\item تولید ویژگی‌های تحلیلی مالی
\item آموزش مدل‌های یادگیری ماشین
\item مدیریت Workflowها توسط Apache Airflow
\item طراحی معماری مبتنی بر Docker Compose
\item پیاده‌سازی API جهت سرویس‌دهی مدل
\item ایجاد زیرساخت اولیه MLOps
\end{itemize}

%===================================================

\section{تعریف مسئله}

در این پروژه قصد داریم از داده‌های بازار بورس ایران جهت توسعه مدل‌های یادگیری ماشین استفاده کنیم.

مدل از داده‌های زیر استفاده خواهد کرد:

\begin{itemize}
\item قیمت باز شدن (Open)
\item قیمت بسته شدن (Close)
\item بیشترین قیمت (High)
\item کمترین قیمت (Low)
\item حجم معاملات (Volume)
\item شاخص‌های تکنیکال
\begin{itemize}
\item RSI
\item MACD
\item Moving Average
\item Lag Features
\end{itemize}
\end{itemize}

خروجی مدل در فاز اولیه می‌تواند شامل موارد زیر باشد:

\begin{enumerate}
\item پیش‌بینی قیمت روز آینده سهام
\item دسته‌بندی تصمیم سرمایه‌گذاری:
\begin{center}
BUY / HOLD / SELL
\end{center}
\end{enumerate}

%===================================================

\section{معماری اولیه سیستم}

معماری پروژه مطابق معماری‌های متداول \lr{MLOps} طراحی شده است.

جریان پردازش سیستم به شکل زیر خواهد بود:

\begin{center}
Iran Stock Data

$\downarrow$

Apache Airflow

$\downarrow$

Data Ingestion

$\downarrow$

Preprocessing

$\downarrow$

Feature Engineering

$\downarrow$

Model Training

$\downarrow$

Model Evaluation

$\downarrow$

MLflow Tracking

$\downarrow$

FastAPI Prediction Service
\end{center}

%===================================================

\section{تکنولوژی‌های استفاده‌شده}

\subsection{Apache Airflow}

استفاده برای:

\begin{itemize}
\item Workflow Orchestration
\item DAG Scheduling
\item Task Management
\end{itemize}

دلیل انتخاب:

Apache Airflow یکی از رایج‌ترین ابزارهای صنعت برای مدیریت Pipelineهای داده و سامانه‌های MLOps محسوب می‌شود.

\subsection{Docker و Docker Compose}

استفاده برای:

\begin{itemize}
\item Containerization
\item Environment Isolation
\item Multi-Service Deployment
\end{itemize}

\subsection{Python}

زبان اصلی توسعه پروژه.

استفاده در:

\begin{itemize}
\item Data Processing
\item Machine Learning
\item API Development
\item Workflow Definition
\end{itemize}

\subsection{PostgreSQL}

استفاده برای:

\begin{itemize}
\item Airflow Metadata Storage
\item Logging
\item Model Metadata
\end{itemize}

\subsection{Redis}

استفاده به عنوان:

\begin{itemize}
\item Message Broker
\item Celery Backend
\end{itemize}

\subsection{Scikit-Learn}

استفاده برای:

\begin{itemize}
\item Model Training
\item Classification
\item Regression
\item Model Evaluation
\end{itemize}

\subsection{MLflow}

استفاده برای:

\begin{itemize}
\item Experiment Tracking
\item Model Registry
\item Metric Logging
\end{itemize}

\subsection{FastAPI}

استفاده برای:

\begin{itemize}
\item REST API Development
\item Model Serving
\item Prediction Endpoints
\end{itemize}

\subsection{Prometheus + Grafana}

(در فازهای آینده پروژه)

برای:

\begin{itemize}
\item Monitoring
\item Metrics Visualization
\end{itemize}

%===================================================

\section{منابع و مستندات مطالعه‌شده}

در فاز اولیه پروژه، اعضای تیم جهت طراحی معماری و انتخاب تکنولوژی‌ها، منابع رسمی زیر را مطالعه کردند.

\begin{longtable}{|p{4cm}|p{9cm}|}

\hline
\textbf{تکنولوژی} & \textbf{منابع مطالعه‌شده} \\
\hline

Apache Airflow &
\url{https://airflow.apache.org/docs/}
\\
\hline

MLflow &
\url{https://mlflow.org/docs/latest/}
\\
\hline

FastAPI &
\url{https://fastapi.tiangolo.com/}
\\
\hline

Redis &
\url{https://redis.io/docs/}
\\
\hline

TSETMC Data Sources &
pytsetmc, tsetmc-api
\\
\hline

\end{longtable}

%===================================================

\section{تقسیم وظایف اعضای تیم}

به‌منظور توسعه ساختاریافته پروژه، مسئولیت‌ها بر اساس تخصص اعضا تقسیم شده است.

\begin{table}[h]

\centering

\begin{tabular}{|p{4cm}|p{8cm}|}

\hline

\textbf{عضو تیم} & \textbf{مسئولیت‌ها} \\

\hline

مهیار قاسمی &
طراحی معماری سیستم، طراحی DAGهای Airflow، زمان‌بندی Workflowها، طراحی وابستگی Taskها، بهینه‌سازی معماری
\\

\hline

محمد حسین صدیقی &
تحلیل داده‌های بورس، Feature Engineering، طراحی اندیکاتورهای مالی، آموزش مدل‌ها، ارزیابی مدل
\\

\hline

سامان زارع &
طراحی Repository Structure، Docker Compose Infrastructure، توسعه FastAPI، Service Integration
\\

\hline

\end{tabular}

\end{table}

%===================================================

\section{نحوه همکاری تیم}

برای مدیریت فرآیند توسعه نرم‌افزار، ساختار همکاری زیر انتخاب شده است:

\begin{itemize}
\item استفاده از GitHub Repository مشترک
\item توسعه مبتنی بر Branch
\item Pull Requests
\item Code Review بین اعضا
\item مستندسازی تغییرات
\end{itemize}

تقسیم کار پروژه به‌صورت ماژولار انجام شده تا امکان توسعه موازی بخش‌های مختلف توسط اعضای تیم فراهم شود.

%===================================================

\section{برنامه توسعه آینده پروژه}

پروژه در هفته‌های آینده توسعه خواهد یافت.

برنامه اولیه شامل:

\begin{enumerate}
\item MLflow Integration
\item Model Registry
\item Data Drift Detection
\item Monitoring Dashboard
\item CI/CD Pipeline
\item Kubernetes Deployment
\end{enumerate}

%===================================================

\section{جمع‌بندی}

در این فاز اولیه، معماری یک سامانه \lr{MLOps} مبتنی بر Apache Airflow جهت تحلیل و پیش‌بینی داده‌های بازار بورس ایران طراحی شد.

در این مرحله:

\begin{itemize}
\item معماری اولیه سیستم مشخص شد.
\item تکنولوژی‌های اصلی انتخاب شدند.
\item منابع مطالعاتی بررسی شدند.
\item تقسیم وظایف بین اعضا انجام شد.
\item زیرساخت اولیه Docker Compose تعریف گردید.
\end{itemize}

این ساختار امکان توسعه تدریجی پروژه در طول ترم را فراهم خواهد کرد.

\end{document}